% Fichier complété par tools/complete_missing_corrections.py.
% Source : COR/COR-04-AP/65_Eclipse_03/65_Eclipse_03.tex
% Statut : rédigé (7 question(s)).
\begin{corrigebox}[Corrigé rédigé]
\CorrigeQuestion{1}
$C_{V2}(p)=(1+k_f\tau_vp)/(1+\tau_vp)$ avec $k_f>1$ est une avance de
            phase. Son zéro est placé avant son pôle; entre les deux, le gain monte de
            $20$ dB/décade et la phase devient positive, ce qui augmente la marge.
\CorrigeQuestion{2}
On lit $\omega_{0dB}$ puis la phase de $W$. Si elle est inférieure à
            $-180^\circ+M_{\varphi}$, la marge est insuffisante; l'avance de phase est
            nécessaire pour remonter la phase autour de cette pulsation.
\CorrigeQuestion{3}
La phase
            $\varphi=\arctan(k_f\tau_v\omega)-\arctan(\tau_v\omega)$ est maximale pour
            $\boxed{\omega_m=1/(\tau_v\sqrt{k_f})}$.
\CorrigeQuestion{4}
On choisit $\omega_m$ égale à la pulsation lue pour la phase cible, donc
            $\boxed{\tau_v=1/(\omega_m\sqrt{k_f})}$.
\CorrigeQuestion{5}
Sur le Black, on translate la courbe verticalement de
            $20\log(K_V/75)$ jusqu'à obtenir la distance maximale au point critique
            $(-180^\circ,0\,dB)$. La valeur correspondante maximise la marge de phase.
\CorrigeQuestion{6}
On mesure pour chaque courbe le premier instant d'entrée définitive dans
            la bande $\pm5\%$. La valeur optimale de $K_V$ est la plus rapide qui respecte
            également le dépassement et la stabilité.
\CorrigeQuestion{7}
L'intégrateur assure le rejet d'une perturbation constante, mais une rampe
            laisse une erreur finie sauf si la boucle possède un intégrateur supplémentaire.
            La robustesse aux rampes n'est donc que partielle.
\end{corrigebox}
